% Partial authored draft only. Not included in review or counted as rewritten.
\chapter{表示距离：统计区分、路径代价与几何模型}
\label{chp:riemannian_base}

前一章讨论了关系基底、局部属性空间和统计区分度。本章进一步追问：系统说两个对象接近，究竟是在比较属性、统计分布、访问路径，还是实现成本？这些比较可以共同参与决策，却不能共享一个未经定义的距离。我们将先规定比较对象与有效域，再分别建立对称距离、统计尺度及路径模型，说明它们何时能够联合，何时必须保持分离。由此形成的几何不是预先给定的全部现实，而是带来源、版本与任务条件的计算模型。

我们还需要区分模型的稳定性和环境事实的约束。目标变化不能自行改变观测证据，事实记录也不能因被称为客观而免于修订。可达路径不等于有效推理，最短路径不等于最正确结论，保范数变换也不自动保证逻辑规则。我们将在固定结构条件下推导路径与局部度量的关系，再把接口更新、群变换和实现预算放回模型。HSF-HD提供这些一般状态动力学条件，AgentOS则在记忆、控制和边界中实例化它们，二者不因共用几何术语而成为同一个理论层次。

\section{比较对象与距离公理：对称性不提供因果方向}
\label{sec:section_203}

我们首先声明距离比较的对象。两个对象记录可以因为颜色接近而具有较小属性距离，也可以因为任务角色不同而具有较大的转换成本。两段程序可以输出相同结果却需要不同运行时间，两条数据库记录可以指向同一对象却具有不同来源。若不明确比较空间，接近一词就会把数值相似、任务等价和执行可达混在一起。我们把对象身份、属性编码、关系位置、证据分布及实现成本分别建模，再通过声明的读出接口使用这些量，而不设想一种距离独立决定对象的全部语义。

在集合 $X$ 上，一个距离函数 $d:X\times X\to[0,\infty)$ 需要非负性、零距离当且仅当对象相同、对称性及三角不等式。若零距离可以对应两个不同对象，我们得到的是伪度量，可以在零距离等价类上再研究商空间。若访问代价具有方向，去往目标和返回起点可能不同，它就不能未经处理成为对称距离。若允许不可达，我们可以使用扩展值或显式不可达标记，但算法必须说明如何处理这些值。任务把两个对象视为等价时，零距离也应说明是在任务商空间中相同，而不是现实身份已经合并。

固定有限维实属性空间中，设坐标按参考尺度归一化，$G$ 为对称正定矩阵。我们定义
\[
 d_G(x,y)=\sqrt{(x-y)^{\mathsf T}G(x-y)}.
\]
由正定性可取可逆矩阵 $C$ 使 $G=C^{\mathsf T}C$，于是距离等于 $\|C(x-y)\|_2$。非负性和对称性直接成立，可逆性使零距离只在两点相同时成立，欧氏范数的三角不等式则给出三角关系。因此我们得到一个真正的距离，而不只是把某个相似评分改名为距离。若矩阵仅半正定，零方向会压缩不同输入，得到的只是伪度量；若矩阵不对称，二次型只看到对称部分，不能借此表达所有方向依赖。

这个构造依赖固定比较规则。若我们对每一对输入临时选择不同的矩阵，对称性和三角不等式就需要重新检查；若矩阵随位置变化，则两点距离通常需要由路径长度的下确界定义，而不是只在一个端点套用固定二次型。在连续几何中，局部正定度量描述小变化的长度，整体距离还依赖空间连通性及允许路径。离散图中的跳数和属性空间中的二次距离也不是天然同一量：前者受边集合限制，后者可能允许连接在图上并不合法的对象。我们只在明确映射及误差条件时联合它们。

对称性尤其不能替代因果或逻辑方向。若某条件是另一事件的必要条件，逻辑上得到的是事件成立蕴含条件成立，不能据此倒转为条件成立就保证事件发生。两个变量统计相关，也不说明干预其中一个必然改变另一个。我们可以在无向层保存相似性，在有向层保存依赖或因果假设，在规则层保存蕴含条件；所有层共享身份与来源，却分别进行验证。一项对称距离小，只说明规定比较下接近，不说明两个命题相互推出，也不说明对象能够互换执行。

火与热的例子可用于说明建模粒度。我们需要先规定变量是某种测量、某类事件标签还是语词出现，并说明采样环境与条件。词语共现所得统计关系和环境测量所得关系不能直接视为同一对象，相关尺度也可能随语境和采样方案改变。把一项关联写入当前模型不会使它成为全宇宙固定距离。模型可以保持对证据的客观约束，同时承认观测不完整及参数可修订；客观性应体现为来源可核查和规则一致，而不是把矩阵数值宣布为不受任何更新影响。

猫和狗在某种哺乳动物属性编码中可能较近，在具体识别任务中却必须保持不同标签。对数函数和香草冰淇淋在一种字段空间中可能缺少直接可比属性，在教学或菜单定价任务中又可以具有有用关系。属性正交不表示图结构一定不连通，数值距离较大也不表示完全没有语义关联。我们让任务提供比较准则、适用范围和读出目标，保留不同准则之间的不一致作为信息，而不强制所有对象进入一个统一远近序列。

结构编辑代价是另一个需要单独定义的对象。若编辑操作及其逆操作具有相同成本，并且操作能连接全部比较对象，最小累计编辑成本可以产生对称路径距离；若操作不允许逆转，或合并与拆分代价不同，通常得到方向代价。零成本编辑还可能使不同记录处于同一伪度量类。我们在接口中列出操作、成本、合法条件和对象等价准则，再分析路径性质。这样，距离公理成为可检查的模型责任，而不是从一种本体论名称自动取得的保证。
